\chapter{Design Técnico e Modelagem Detalhada}
\label{chap:design_tecnico}

\section{Mapeamento Objeto-Relacional (ORM)}
O mapeamento objeto-relacional traduz o modelo conceitual de domínio para esquemas físicos de persistência em banco relacional, garantindo integridade referencial através de chaves primárias e estrangeiras bem definidas.

\subsection{Esquema Físico Textual}
\begin{lstlisting}[language=SQL, caption={Esquema Relacional das Tabelas Principais (DDL Parcial)}]
CREATE TABLE tb_usuario (
    id UUID PRIMARY KEY DEFAULT gen_random_uuid(),
    nome VARCHAR(120) NOT NULL,
    email VARCHAR(120) UNIQUE NOT NULL,
    senha_hash VARCHAR(255) NOT NULL,
    papel VARCHAR(30) NOT NULL,
    ativo BOOLEAN DEFAULT TRUE,
    criado_em TIMESTAMP WITH TIME ZONE DEFAULT CURRENT_TIMESTAMP
);

CREATE TABLE tb_agendamento (
    id UUID PRIMARY KEY DEFAULT gen_random_uuid(),
    cliente_id UUID NOT NULL REFERENCES tb_usuario(id),
    profissional_id UUID NOT NULL REFERENCES tb_usuario(id),
    data_hora TIMESTAMP WITH TIME ZONE NOT NULL,
    status VARCHAR(20) DEFAULT 'AGENDADO',
    CONSTRAINT unq_horario_profissional UNIQUE (profissional_id, data_hora)
);
\end{lstlisting}

\section{Diagrama de Classes em Nível de Projeto}
Diferente do diagrama de análise, o diagrama de projeto especifica as classes de implementação: controladores, casos de uso, repositórios, entidades e DTOs, incluindo visibilidade (\texttt{+}, \texttt{-}, \texttt{\#}), tipos de retorno e parâmetros de métodos. A \autoref{fig:cls_projeto_agendamento} ilustra o diagrama de classes detalhado do módulo de agendamento.

\paragraph{Diretriz de Decomposição de Diagramas:}
Para evitar poluição visual e cruzamento desordenado de linhas, o diagrama de classes em nível de projeto deve ser apresentado dividido por módulos funcionais (subpacotes), acompanhado de tabela sintética de métodos e atributos.

% PLACEHOLDER STARUML v7.0: Diagrama de Classes de Projeto
% Ferramenta: StarUML v7.0 -> Model -> Add Diagram -> Class Diagram
% Arquivo: Imagens/cls_projeto_agendamento.png (Exportar em PNG 300 DPI ou PDF vetorial)
% Descomente a linha abaixo apos salvar o arquivo no diretorio Imagens/:
% \incluirdiagrama{Imagens/cls_projeto_agendamento.png}{Diagrama de Classes de Projeto -- Módulo Agendamento}{fig:cls_projeto_agendamento}
\begin{figure}[htbp]
    \centering
    \fbox{\parbox{0.9\textwidth}{\centering\vspace{1cm}%
        \textbf{[Diagrama de Classes de Projeto -- StarUML v7.0]}\\%
        \textit{Modelar no StarUML v7.0 e exportar para: \texttt{Imagens/cls\_projeto\_agendamento.png}}\\%
        \small Menu: \texttt{File -> Export Diagram as -> PNG...} (Fundo branco, alta resolução 300 DPI)\\%
        \vspace{0.3cm}%
        \footnotesize Para ativar a imagem real, descomente no arquivo \texttt{.tex}:\\
        \texttt{\textbackslash incluirdiagrama\{Imagens/cls\_projeto\_agendamento.png\}\{Classes de Projeto -- Agendamento\}\{fig:cls\_projeto\_agendamento\}}%
    \vspace{1cm}}}
    \caption{Diagrama de Classes de Projeto do Módulo Agendamento (Placeholder StarUML v7.0)}
    \label{fig:cls_projeto_agendamento}
\end{figure}

\section{Diagramas Dinâmicos}

\subsection{Diagramas de Sequência}
Representam a ordem temporal das trocas de mensagens entre objetos e serviços durante a execução de um caso de uso crítico. A \autoref{fig:seq_agendamento} detalha a interação para o caso de uso UC01.

% PLACEHOLDER STARUML v7.0: Diagrama de Sequência
% Ferramenta: StarUML v7.0 -> Model -> Add Diagram -> Sequence Diagram
% Arquivo: Imagens/seq_agendamento.png (Exportar em PNG 300 DPI ou PDF vetorial)
% Descomente a linha abaixo apos salvar o arquivo no diretorio Imagens/:
% \incluirdiagrama{Imagens/seq_agendamento.png}{Diagrama de Sequência -- UC01 Realizar Agendamento}{fig:seq_agendamento}
\begin{figure}[htbp]
    \centering
    \fbox{\parbox{0.9\textwidth}{\centering\vspace{1cm}%
        \textbf{[Diagrama de Sequência -- UC01 Agendamento -- StarUML v7.0]}\\%
        \textit{Modelar no StarUML v7.0 e exportar para: \texttt{Imagens/seq\_agendamento.png}}\\%
        \small Menu: \texttt{File -> Export Diagram as -> PNG...} (Fundo branco, alta resolução 300 DPI)\\%
        \vspace{0.3cm}%
        \footnotesize Para ativar a imagem real, descomente no arquivo \texttt{.tex}:\\
        \texttt{\textbackslash incluirdiagrama\{Imagens/seq\_agendamento.png\}\{Diagrama de Sequência -- UC01\}\{fig:seq\_agendamento\}}%
    \vspace{1cm}}}
    \caption{Diagrama de Sequência do UC01 -- Agendamento (Placeholder StarUML v7.0)}
    \label{fig:seq_agendamento}
\end{figure}

\subsection{Diagramas de Máquinas de Estados}
Modelam o ciclo de vida e as transições de estado permitidas para entidades centrais da regra de negócio (ex.: Consulta, Pedido, Prescrição). A \autoref{fig:dsm_agendamento} demonstra os estados possíveis da entidade Agendamento.

% PLACEHOLDER STARUML v7.0: Diagrama de Máquinas de Estados
% Ferramenta: StarUML v7.0 -> Model -> Add Diagram -> Statechart Diagram
% Arquivo: Imagens/dsm_agendamento.png (Exportar em PNG 300 DPI ou PDF vetorial)
% Descomente a linha abaixo apos salvar o arquivo no diretorio Imagens/:
% \incluirdiagrama{Imagens/dsm_agendamento.png}{Máquina de Estados da Entidade Agendamento}{fig:dsm_agendamento}
\begin{figure}[htbp]
    \centering
    \fbox{\parbox{0.9\textwidth}{\centering\vspace{1cm}%
        \textbf{[Diagrama de Máquina de Estados -- Entidade Agendamento -- StarUML v7.0]}\\%
        \textit{Modelar no StarUML v7.0 e exportar para: \texttt{Imagens/dsm\_agendamento.png}}\\%
        \small Menu: \texttt{File -> Export Diagram as -> PNG...} (Fundo branco, alta resolução 300 DPI)\\%
        \vspace{0.3cm}%
        \footnotesize Para ativar a imagem real, descomente no arquivo \texttt{.tex}:\\
        \texttt{\textbackslash incluirdiagrama\{Imagens/dsm\_agendamento.png\}\{Máquina de Estados -- Agendamento\}\{fig:dsm\_agendamento\}}%
    \vspace{1cm}}}
    \caption{Diagrama de Máquina de Estados da Entidade Agendamento (Placeholder StarUML v7.0)}
    \label{fig:dsm_agendamento}
\end{figure}

\begin{table}[htbp]
\caption{Matriz de Transição de Estados da Entidade Agendamento}
\label{tab:transicao_estados}
\centering
\small
\begin{tabularx}{\textwidth}{@{} l Y Y l @{}}
\toprule
\textbf{Estado Origem} & \textbf{Evento / Ação} & \textbf{Condição de Guarda} & \textbf{Estado Destino} \\
\midrule
Solicitado & Confirmar Agendamento & Horário válido e sem colisão & Agendado \\
Agendado & Iniciar Atendimento & Presença do cliente confirmada & Em Atendimento \\
Agendado & Cancelar & Solicitação com antecedência $\ge 2$h & Cancelado \\
Em Atendimento & Finalizar Consulta & Prontuário preenchido e assinado & Concluído \\
\bottomrule
\end{tabularx}
\end{table}

\subsection{Diagramas de Atividades}
Ilustram fluxos operacionais completos com bifurcações condicionais, paralelismo (fork/join) e raias de responsabilidade (swimlanes) entre diferentes setores ou sistemas. A \autoref{fig:act_agendamento} mapeia as raias e decisões do processo de agendamento.

% PLACEHOLDER STARUML v7.0: Diagrama de Atividades
% Ferramenta: StarUML v7.0 -> Model -> Add Diagram -> Activity Diagram
% Arquivo: Imagens/act_fluxo_agendamento.png (Exportar em PNG 300 DPI ou PDF vetorial)
% Descomente a linha abaixo apos salvar o arquivo no diretorio Imagens/:
% \incluirdiagrama{Imagens/act_fluxo_agendamento.png}{Fluxo Operacional de Atividades do Agendamento}{fig:act_agendamento}
\begin{figure}[htbp]
    \centering
    \fbox{\parbox{0.9\textwidth}{\centering\vspace{1cm}%
        \textbf{[Diagrama de Atividades com Raias -- StarUML v7.0]}\\%
        \textit{Modelar no StarUML v7.0 e exportar para: \texttt{Imagens/act\_fluxo\_agendamento.png}}\\%
        \small Menu: \texttt{File -> Export Diagram as -> PNG...} (Fundo branco, alta resolução 300 DPI)\\%
        \vspace{0.3cm}%
        \footnotesize Para ativar a imagem real, descomente no arquivo \texttt{.tex}:\\
        \texttt{\textbackslash incluirdiagrama\{Imagens/act\_fluxo\_agendamento.png\}\{Diagrama de Atividades\}\{fig:act\_agendamento\}}%
    \vspace{1cm}}}
    \caption{Diagrama de Atividades do Fluxo de Agendamento (Placeholder StarUML v7.0)}
    \label{fig:act_agendamento}
\end{figure}

\section{Status dos Casos de Uso Implementáveis}
\begin{table}[htbp]
\caption{Quadro de Status e Cobertura de Implementação}
\label{tab:status_implementacao}
\centering
\small
\begin{tabularx}{\textwidth}{@{} l Y c l @{}}
\toprule
\textbf{Caso de Uso} & \textbf{Nome do Módulo} & \textbf{Status Atual} & \textbf{Responsável} \\
\midrule
UC01 & Autenticação e RBAC & \badgeconcluido & Back-end Team \\
UC02 & Agendamento e Grade & \badgeconcluido & Back-end Team \\
UC03 & Prontuário Eletrônico & \badgeemprogresso & Full-stack Team \\
UC04 & Gestão de Estoque & \badgenaoiciada & Back-end Team \\
\bottomrule
\end{tabularx}
\end{table}
